\documentclass[11pt]{letter}

\usepackage[margin=0.75in]{geometry}
\usepackage{setspace}
\usepackage{hyperref}
\usepackage{xcolor}
\usepackage{enumitem}
\usepackage{parskip}

\hypersetup{
    colorlinks=true,
    urlcolor=black,
    linkcolor=black
}

\setlength{\parindent}{0pt}
\setstretch{1.05}
\date{}
\begin{document}

\begin{letter}{}

\opening{Dear Hiring Committee,}

I saw the job post on the \textbf{Salesforce careers website} and thought I'd be a good fit for it.

I am a Master's student in Computer Science at the University of Cincinnati, graduating December 2026, looking for a role as a \textbf{Software Engineering AMTS (College Grad)} at \textbf{Salesforce}. With over two years of production software engineering experience at Siemens Industry Software, I have built backend automation, full stack applications, and cloud platforms. I am excited about the opportunity to contribute to \textbf{Salesforce} and help build scalable, trusted cloud products for customers.

\textbf{Here are reasons that make me a great fit for the role:}

\begin{itemize}[leftmargin=*, itemsep=0.45em]
    \item \textbf{Building scalable products:} At Siemens, I migrated 40+ production environments to Kubernetes using ArgoCD, Helm, and GitLab CI/CD, executed 44 production environment upgrades in 4 months, and designed an Ansible based Blue-Green deployment framework that reduced upgrade downtime by 50\%.

    \item \textbf{Test strategy and automation:} As Quality Coach and SonarQube Specialist, I increased automated code coverage from 35\% to 78\%, and I developed a Go and Terratest framework that caught 95\% of misconfigurations before production. In my Spendora project, I wrote 80 tests across unit, database, and end to end layers.

    \item \textbf{Professional coding in Python, JavaScript, SQL, and Java:} I automated Teamcenter X workflows with Python and Boto3, cutting deployment time by 40\%, and built a React and MySQL application that cut document retrieval time by 90\%. My Computer Engineering degree included object oriented programming coursework.

    \item \textbf{Working thoughtfully with LLMs and agentic tools:} In my graduate research, I built a memory augmented LLM pipeline with Gemini that achieved a 25.7x latency reduction over the ViBR baseline. In PolicyPilot, I designed a 17 tool MCP server for external AI agents and benchmarked hallucination rates before choosing Gemini.

    \item \textbf{Metrics and collaboration:} I used CloudWatch utilization metrics to help cut monthly AWS spend by 40\% (\$15K to \$9K), and I mentored 5 software engineers through structured onboarding.
\end{itemize}

I am available for full time work after my graduation in \textbf{December 2026}, which falls within the Q4 and Q1 start window for this role. As a Master's student at the University of Cincinnati and someone who values professionalism and growth, I am enthusiastic about contributing to \textbf{Salesforce} and its software engineering teams.

Please find my CV attached below. I look forward to hearing from you. \textbf{Salesforce's vision of humans and AI agents driving customer success together is an exciting direction, and I would love to be part of it}.

{\setlength{\parskip}{0pt}
Regards,\\
Tanmay Bhuskute\\
bhuskutv@mail.uc.edu $\mid$ tanmay.v.bhuskute@gmail.com $\mid$ +1 (732) 858-3027\\
3424 Brookline Ave, Flat 19, Cincinnati, OH 45220
}

\end{letter}

\end{document}
